\documentclass[10pt,a4paper]{amsart}
\usepackage[margin=19mm]{geometry}
\usepackage{amsmath,amssymb,mathtools}
\usepackage[scr=rsfs,cal=euler]{mathalfa}
\usepackage{xfrac}
\usepackage[unicode,hidelinks,bookmarks=false]{hyperref}
\newcommand{\ie}{i.e.\ }
\newcommand{\cf}{cf.\ }
\newcommand{\resp}{respectively\ }
\newcommand{\Subsection}{\subsection}
\providecommand{\nobreakdash}{\nobreak\hbox{-}}
\setlength{\emergencystretch}{2em}
\multlinegap0pt
\usepackage{amssymb}
\newtheorem{thm}{Theorem}
\title{Multiplier modules of Hilbert C*-modules revisited}
\author{Michael Frank}
\date{Source: arXiv:2502.17959v3 (CC BY 4.0)}
\begin{document}
\maketitle

\noindent\emph{Source and licence.} Multiplier modules of Hilbert C*-modules revisited. Version arXiv:2502.17959v3 (CC BY 4.0), distributed by arXiv under the Creative Commons Attribution 4.0 International licence, http://creativecommons.org/licenses/by/4.0/, as recorded in the arXiv licence metadata for this exact version and shown on the abstract page https://arxiv.org/abs/2502.17959v3. Full licence notice: \texttt{licenses/arxiv-2502.17959-CC-BY-4.0.txt}.\par\smallskip

\noindent\textit{Editorial scope.} Counted unit: the article's thm funct (source0.tex lines 298--304). The article's complete original proof of it is delivered with it (source0.tex lines 306--310, one \texttt{proof} environment of 5 lines). No mathematics is written, repaired or re-derived: the statement and the proof are the article's own lines, in the article's own notation, and no retained line is substituted or re-flowed.  No further source-local context is needed: every cross-reference of every retained line resolves inside the two delivered spans.  Nothing was omitted from any retained span, so the package carries no omission record.  No retained line contains a citation, so the wrapper carries no bibliography and no bibliography entry.  No retained line contains a figure environment, an image-inclusion command or drawing code, so no image is delivered and the package declares no asset dependency.  Editorial formatting only, in addition: an AMS \texttt{amsart} preamble that restates the article's own declarations and macros used by the retained lines, namely the environments \texttt{thm} on separate counters, together with the standard packages those lines need.  The retained environments print in this document as Theorem 1 in order of appearance, whereas the article prints the same items with the numbering of its own full document.  The complete native source of the article as distributed in the arXiv e-print for this exact version is bundled unchanged as \texttt{sources/173821/source0.tex}.
\begin{thm} \label{funct}
Let $A$ be a C*-algebra with multiplier algebra $M(A)$. Let $X$ be a full Hilbert $A$-module and $M(X)$ be its full multiplier module.
\begin{enumerate}
  \item There does not exist any bounded $M(A)$-linear map $f_0: M(X) \to M(A)$ such that $f_0 \neq 0$ on $M(X)$, but $f_0 =0$ on $X \subseteq M(X)$.
  \item The Banach $M(A)$-module $M(X)'_{M(A)}$ admits an isometric modular embedding into the Banach $A$-module $X'_A$ by restricting an element on the domain from $M(X)$ to $X\subseteq M(X)$. There exist examples such that $X'_A$ is strictly larger than the embedded copy of $M(X)'_{M(A)}$.
\end{enumerate}
\end{thm}

\begin{proof}
Suppose, there exists a bounded $M(A)$-linear functional $f_0 : M(X) \to M(A)$ such that $f_0=0$ on $X \subseteq M(X)$, but $f_0(m) \neq 0$ for some $m \in M(X)$. Let $\{ x_\alpha: \alpha \in I\}$ be a net of elements of $X$ converging strictly to $m$, i.e.~the nets $\{ x_\alpha a: \alpha \in I\}$ converge to $ma \in X$ in norm for any $a \in A$. Consider the set $\{ f_0(ma) : a \in A \}$. All these values are equal to zero by supposition. Since  $f_0(ma) = f_0(m) a = 0$ for any $a \in A$ and $A$ is an essential ideal of $M(A)$ we conclude $f_0(m) =0$, a contradiction to our assumption.

Restricting $f \in M(X)'$ to $X \subseteq M(X)$ we obtain a bounded $A$-linear functional of $X'$. The norm is preserved, since $X$ ist strictly dense in $M(X)$ and there does not exist any non-trivial bounded $M(A)$-linear functional on $M(X)$ vanishing on $X \subseteq M(X)$. The example in the beginning of the present section can be read as follows: Let $A$ be a C*-algebra such that $LMA) \supset M(A)$. Setting $A=X$ we get $X'=LM(A)$ and $M(X)'=M(A)'= M(A)$ since $M(X)=M(X)'$ is selfdual. So $X' \supset M(X)'$ and an $A$-valued bounded functional in $LM(A) \setminus M(A)$ cannot be continued from $X \subseteq M(X)$ to $M(X)$.
\end{proof}

\end{document}
